\chapter{माप सिद्धांत}\label{ch:b3:measure}

$\R$ का कोई उपसमुच्चय कितना लंबा होता है? भोला उत्तर --- प्रत्येक
समुच्चय को अंतरालों की लंबाई का विस्तार करने वाली कोई
स्थानांतरण-अपरिवर्ती लंबाई सौंप देना --- \emph{असंभव} है: इस
अध्याय के अंत में विटाली की रचना एक ऐसा समुच्चय बना देती है जिसकी
कोई संगत लंबाई नहीं। \omterm{def:b3:measure:measure}{माप} सिद्धांत अनुशासित पीछे हटना है: हम अपना
ध्यान \emph{मापनीय} समुच्चयों के एक समृद्ध वर्ग तक सीमित कर लेते
हैं, जिस पर गणनीय रूप से योज्य लंबाई विद्यमान है और अद्वितीय भी।
पुरस्कार अपार हैं --- लेबेग का समाकल (\cref{ch:b3:lebesgue}), कार्यकीय विश्लेषण की
$L^p$ समष्टियाँ, और समूची आधुनिक प्रायिकता (\cref{ch:b3:probability}) यहाँ सिद्ध होने
वाली तीन प्रमेयों पर खड़ी हैं: डिन्किन की अद्वितीयता प्रमेयिका,
कैराथियोडोरी की विस्तार प्रमेय, और \omterm{def:b3:measure:lebesgueouter}{लेबेग माप} का अस्तित्व।

\section{$\sigma$-बीजगणित}

\begin{definition}\label{def:b3:measure:sigmaalgebra}
किसी समुच्चय $X$ पर \emph{$\sigma$-बीजगणित}\index{सिग्मा-बीजगणित@$\sigma$-बीजगणित}
उपसमुच्चयों का ऐसा कुल $\mathcal A$ है जिसमें $\varnothing$ हो और जो पूरक तथा
\emph{गणनीय} सम्मिलनों के अंतर्गत स्थिर हो (और इसलिए गणनीय
प्रतिच्छेदनों, समुच्चय अंतरों के अंतर्गत भी, तथा जिसमें $X$ हो)।
युग्म $(X, \mathcal A)$ \emph{मापनीय समष्टि} कहलाता है; $\mathcal A$ के सदस्य
\emph{मापनीय समुच्चय} कहलाते हैं। उपसमुच्चयों के किसी भी कुल $\mathcal E$
के लिए $\sigma(\mathcal E)$ $\mathcal E$ को अंतर्विष्ट करने वाला सबसे छोटा
$\sigma$-बीजगणित निरूपित करता है (उन सबका प्रतिच्छेदन ---
$\sigma$-बीजगणितों का प्रतिच्छेदन स्वयं वैसा ही होता है)।
\end{definition}

\begin{definition}\label{def:b3:measure:borel}
किसी संस्थितिक समष्टि का \emph{बोरेल $\sigma$-बीजगणित}\index{बोरेल
सिग्मा-बीजगणित@%
बोरेल $\sigma$-बीजगणित} $\mathcal B(X) = \sigma(\{\text{विवृत समुच्चय}\})$ है। $\R$ पर: $\mathcal B(\R)$ विवृत अंतरालों से,
संवृत अंतरालों से, किरणों $(-\infty, a]$ से, और परिमेय सिरों वाली किरणों से
भी जनित होता है (\cref{exo:b3:measure:1}) --- प्रत्येक कुल गणनीय संक्रियाओं से \omterm{def:b3:topology:topology}{विवृत
समुच्चय} जनित कर देता है, उदाहरणार्थ $\R$ का प्रत्येक \omterm{def:b3:topology:topology}{विवृत
समुच्चय} परिमेय आँकड़ों वाले विवृत अंतरालों का गणनीय सम्मिलन होता
है।
\end{definition}

\begin{definition}\label{def:b3:measure:dynkin}
\emph{$\pi$-प्रणाली}\index{पाई-प्रणाली@$\pi$-प्रणाली} ऐसा कुल है जो
परिमित प्रतिच्छेदनों के अंतर्गत स्थिर हो।
\emph{$\lambda$-प्रणाली}\index{लैम्ब्डा-प्रणाली@$\lambda$-प्रणाली}
(डिन्किन वर्ग) ऐसा कुल $\mathcal D$ है जिसके लिए: $X \in \mathcal
D$; $A, B \in \mathcal D$, $A \subseteq B$ $\Rightarrow$
$B
\setminus A \in \mathcal D$; और $A_n \uparrow A$, $A_n \in
\mathcal D$ $\Rightarrow$ $A \in \mathcal D$।
\end{definition}

\begin{theorem}[डिन्किन की $\pi$--$\lambda$ प्रमेयिका]
\label{thm:b3:measure:dynkin}
यदि कोई $\lambda$-प्रणाली $\mathcal D$ किसी $\pi$-प्रणाली $\mathcal P$ को अंतर्विष्ट
करे, तो $\mathcal D \supseteq
\sigma(\mathcal P)$।\index{डिन्किन की प्रमेयिका}
\end{theorem}

\begin{proof}
मान लीजिए $\mathcal D_0$ $\mathcal P$ को अंतर्विष्ट करने वाली सबसे छोटी
$\lambda$-प्रणाली है (ऐसी सबका प्रतिच्छेदन); यह दिखाना पर्याप्त है कि
$\mathcal D_0$ एक $\sigma$-बीजगणित है, क्योंकि तब $\sigma(\mathcal P) \subseteq \mathcal D_0 \subseteq \mathcal
D$। परिमित प्रतिच्छेदनों के
अंतर्गत स्थिर कोई $\lambda$-प्रणाली \emph{होती} ही $\sigma$-बीजगणित है:
पूरक ($A \subseteq X$ के साथ $X \setminus A = X
\setminus A$), परिमित सम्मिलन ($A \cup B =
X \setminus ((X\setminus A)\cap(X\setminus B))$), और $\bigcup_{k \leq n}A_k \uparrow \bigcup_kA_k$ के
माध्यम से गणनीय सम्मिलन। अतः हम दो चरणों में सिद्ध करते हैं कि
$\mathcal D_0$ एक $\pi$-प्रणाली है। मान लीजिए
\[
\mathcal D_1 = \{A \in \mathcal D_0 : A \cap P \in \mathcal
D_0 \ \forall P \in \mathcal P\}.
\]
$\mathcal D_1$ एक $\lambda$-प्रणाली है (तीनों स्वयंसिद्ध $P$ से प्रतिच्छेदन
लेकर सत्यापित हो जाते हैं: उदाहरणार्थ $(B\setminus A)\cap P
= (B \cap P)\setminus(A \cap P)$, जो $\mathcal D_0$ के भीतर एक
उचित अंतर है) और उसमें $\mathcal P$ ($\pi$-प्रणाली) है: अतः $\mathcal D_1 = \mathcal D_0$। अब मान
लीजिए
\[
\mathcal D_2 = \{A \in \mathcal D_0 : A \cap D \in \mathcal
D_0\ \forall D \in \mathcal D_0\}.
\]
पिछले चरण से $\mathcal D_2 \supseteq \mathcal P$; और उसी सत्यापन से $\mathcal D_2$ एक $\lambda$-प्रणाली है:
अतः $\mathcal D_2 = \mathcal D_0$, जिसका ठीक-ठीक अर्थ है कि $\mathcal D_0$ प्रतिच्छेदनों के अंतर्गत
स्थिर है।
\end{proof}

\section{माप}

\begin{definition}\label{def:b3:measure:measure}
$(X, \mathcal A)$ पर \emph{माप}\index{माप} ऐसा प्रतिचित्रण $\mu \colon \mathcal A \to [0, +\infty]$ है जिसके लिए
$\mu(\varnothing)
= 0$ हो और जो \emph{$\sigma$-योज्य} हो: जोड़े-जोड़े में असंयुक्त
$(A_n)_{n\in\N}$ के लिए
\[
\mu\Bigl(\bigsqcup_n A_n\Bigr) = \sum_n \mu(A_n).
\]
$(X, \mathcal A, \mu)$ एक \emph{माप समष्टि} है; $\mu$ \emph{परिमित} कहलाता है यदि
$\mu(X) < \infty$ हो, \emph{प्रायिकता माप} यदि $\mu(X) = 1$ हो, और \emph{$\sigma$-परिमित}
यदि $X$ परिमित माप वाले समुच्चयों का गणनीय सम्मिलन हो। उदाहरण:
$(\N, \mathcal P(\N))$ पर गणना माप; डिराक द्रव्यमान $\delta_a(A)
= \mathbf 1_{a \in A}$; और, इस अध्याय की मुख्य
वस्तु, \omterm{def:b3:measure:lebesgueouter}{लेबेग माप}।
\end{definition}

\begin{proposition}\label{prop:b3:measure:basics}
मान लीजिए $\mu$ एक \omterm{def:b3:measure:measure}{माप} है। (क) एकदिष्टता: $A \subseteq B
\Rightarrow \mu(A) \leq \mu(B)$। (ख) गणनीय
उपयोज्यता: $\mu(\bigcup A_n) \leq \sum\mu(A_n)$। (ग) नीचे से \omterm{def:b3:topology:continuity}{सांतत्य}: $A_n \uparrow A \Rightarrow \mu(A_n) \to \mu(A)$। (घ) ऊपर से \omterm{def:b3:topology:continuity}{सांतत्य}:
$A_n \downarrow A$ \emph{जिसमें $\mu(A_1) < \infty$} $\Rightarrow \mu(A_n) \to \mu(A)$।
\end{proposition}

\begin{proof}
(क) $B = A \sqcup (B\setminus A)$। (ख) असंयुक्त बनाइए: $B_n = A_n
\setminus \bigcup_{k<n}A_k$ असंयुक्त हैं और उनका
सम्मिलन वही है, तथा $\mu(B_n) \leq \mu(A_n)$। (ग) $A = \bigsqcup_n (A_n
\setminus A_{n-1})$ ($A_0 = \varnothing$): $\sum\mu(A_n\setminus A_{n-1})$ के आंशिक
योग $\mu(A_n)$ हैं। (घ) $A_1 \setminus A_n \uparrow A_1 \setminus A$ पर (ग) लगाइए और $\mu(A_1)$ में से घटाइए ---
परिमितता ही इस घटाव को वैध बनाती है। उसके बिना प्रतिउदाहरण: \omterm{def:b3:measure:lebesgueouter}{लेबेग
माप} के लिए $A_n = [n, \infty)$: $A_n \downarrow \varnothing$ पर $\mu(A_n) =
\infty$।
\end{proof}

\begin{theorem}[अद्वितीयता]\label{thm:b3:measure:uniqueness}
मान लीजिए $\mu, \nu$ $\sigma(\mathcal P)$ पर \omterm{def:b3:measure:measure}{माप} हैं, $\mathcal
P$ एक $\pi$-प्रणाली है,
और $\mathcal P$ पर $\mu = \nu$। यदि $P_k \uparrow X$ और $\mu(P_k) < \infty$ वाले समुच्चय $P_k \in \mathcal P$ हों,
तो समूचे $\sigma(\mathcal P)$ पर $\mu = \nu$।
\end{theorem}

\begin{proof}
$k$ स्थिर कीजिए और $\sigma(\mathcal
P)$ पर परिमित मापों $\mu_k(A) = \mu(A \cap
P_k)$ और $\nu_k(A) = \nu(A \cap P_k)$ पर विचार
कीजिए: वे $\mathcal P$ पर सहमत हैं, क्योंकि $P \cap P_k \in \mathcal
P$ ($\pi$-प्रणाली), और वे
$X$ को वही परिमित मान $\mu(P_k)$ सौंपते हैं। वर्ग $\mathcal D = \{A : \mu_k(A) =
\nu_k(A)\}$ एक
$\lambda$-प्रणाली है: $X \in \mathcal D$; उचित अंतर घटाव से (परिमित मान); और बढ़ती
सीमाएँ नीचे से \omterm{def:b3:topology:continuity}{सांतत्य} से (\cref{prop:b3:measure:basics}(ग))। उसमें $\pi$-प्रणाली $\mathcal P$
अंतर्विष्ट है, अतः डिन्किन (\cref{thm:b3:measure:dynkin}) से $\mathcal D \supseteq
\sigma(\mathcal P)$: इसलिए सर्वत्र $\mu_k = \nu_k$।
अंततः, किसी भी $A \in \sigma(\mathcal P)$ के लिए $A
\cap P_k \uparrow A$ के अनुदिश नीचे से \omterm{def:b3:topology:continuity}{सांतत्य}
$\mu(A) = \lim_k\mu_k(A) =
\lim_k\nu_k(A) = \nu(A)$ दे देती है।
\end{proof}

\section{बाह्य माप और कैराथियोडोरी की प्रमेय}

\begin{definition}\label{def:b3:measure:outer}
$X$ पर \emph{बाह्य माप}\index{बाह्य माप} ऐसा प्रतिचित्रण $\mu^* \colon \mathcal P(X) \to [0, \infty]$
है जिसके लिए $\mu^*(\varnothing) = 0$ हो, जो एकदिष्ट हो और गणनीय रूप से उपयोज्य।
कोई समुच्चय $A$ \emph{$\mu^*$-मापनीय} (कैराथियोडोरी) कहलाता है
यदि वह प्रत्येक समुच्चय को योज्य ढंग से बाँट दे:
\[
\mu^*(E) = \mu^*(E \cap A) + \mu^*(E \setminus A)
\qquad \text{प्रत्येक } E \subseteq X \text{ के लिए}
\]
($\leq$ उपयोज्यता से सदैव लागू होता है; सार $\geq$ में है)।
\end{definition}

\begin{theorem}[कैराथियोडोरी]\label{thm:b3:measure:caratheodory}
$\mu^*$-मापनीय समुच्चय एक $\sigma$-बीजगणित $\mathcal
M$ बनाते हैं, और $\mu^*\restriction_{\mathcal M}$
एक \omterm{def:b3:measure:measure}{माप} है। इसके अतिरिक्त $\mu^*(N) = 0$ वाला प्रत्येक समुच्चय $\mathcal M$ में
होता है (\omterm{def:b3:measure:measure}{माप} \emph{\omterm{def:b3:complete:complete}{पूर्ण}} है)।\index{कैराथियोडोरी की प्रमेय}
\end{theorem}

\begin{proof}
$\mathcal M$ में $\varnothing$ है और वह पूरक के अंतर्गत स्थिर है (परिभाषित शर्त
$A$, $X\setminus A$ में सममित है)। \emph{परिमित सम्मिलन}: मान लीजिए
$A, B \in \mathcal
M$ और $E$ मनमाना है; $E$ को $A$ से बाँटिए, फिर प्रत्येक
टुकड़े को $B$ से:
\[
\mu^*(E) = \mu^*(E\cap A\cap B) + \mu^*(E\cap A\setminus B) +
\mu^*(E\cap B\setminus A) + \mu^*(E\setminus(A\cup B)).
\]
पहले तीन टुकड़े $E \cap (A \cup B)$ को आच्छादित करते हैं, अतः उपयोज्यता से
$\mu^*(E) \geq \mu^*(E\cap(A\cup B)) +
\mu^*(E\setminus(A\cup B))$ मिलता है: इसलिए $A \cup B \in \mathcal M$। आगमन से परिमित सम्मिलन; और पूरकों के
साथ परिमित असंयुक्तता संबंधी हेरफेर उपलब्ध हो जाते हैं।

\emph{$\mathcal M$ पर योज्यता}: असंयुक्त $A, B \in
\mathcal M$ और किसी भी $E$ के लिए
$\mu^*(E\cap(A\sqcup B)) = \mu^*(E\cap
A) + \mu^*(E \cap B)$ ($A$ से बाँटिए); और आगमन से
\[
\mu^*\Bigl(E \cap \bigsqcup_{k\leq n}A_k\Bigr)
= \sum_{k\leq n}\mu^*(E\cap A_k).
\tag{$*$}
\]

\emph{गणनीय सम्मिलन}: मान लीजिए $(A_k) \subseteq \mathcal M$ असंयुक्त हैं (बीजगणित
$\mathcal M$ के भीतर असंयुक्त बनाने से यह पर्याप्त है), $A = \bigsqcup A_k$, और $E$
मनमाना। $\bigsqcup_{k\leq n}A_k \in \mathcal M$ तथा एकदिष्टता का उपयोग करते हुए:
\[
\mu^*(E) = \mu^*\Bigl(E\cap\bigsqcup_{k\leq n}A_k\Bigr) +
\mu^*\Bigl(E\setminus\bigsqcup_{k\leq n}A_k\Bigr)
\geq \sum_{k \leq n}\mu^*(E\cap A_k) + \mu^*(E\setminus A)
\]
($*$) से। अब $n \to \infty$ लीजिए और गणनीय उपयोज्यता उलटकर लगाइए:
\[
\mu^*(E) \geq \sum_{k}\mu^*(E\cap A_k) + \mu^*(E\setminus A)
\geq \mu^*(E \cap A) + \mu^*(E\setminus A) \geq \mu^*(E):
\]
अतः सभी असमिकाएँ समताएँ हैं। इससे $A \in
\mathcal M$ भी सिद्ध होता है और,
$E = A$ लेने पर, $\mathcal M$ पर $\mu^*$ की गणनीय योज्यता भी।

\emph{शून्य समुच्चय}: यदि $\mu^*(N) = 0$ हो, तो किसी भी $E$ के लिए
$\mu^*(E \cap N) + \mu^*(E\setminus N) \leq 0 + \mu^*(E)$: अतः $N \in \mathcal M$।
\end{proof}

\section{\texorpdfstring{$\R$}{R} पर लेबेग माप}

\begin{definition}\label{def:b3:measure:lebesgueouter}
$A
\subseteq \R$ का \emph{लेबेग बाह्य माप}\index{लेबेग माप} है
\[
\lambda^*(A) = \inf\Bigl\{\sum_{n} (b_n - a_n) :
A \subseteq \bigcup_n \intoo{a_n}{b_n}\Bigr\}
\]
(विवृत अंतरालों से गणनीय आवरण)।
\end{definition}

\begin{lemma}\label{lem:b3:measure:outerbasics}
$\lambda^*$ एक \omterm{def:b3:measure:outer}{बाह्य माप} है, जो स्थानांतरणों के अंतर्गत अपरिवर्ती है,
और प्रत्येक अंतराल $I$ के लिए $\lambda^*(I) = \ell(I)$ (उसकी लंबाई)।
\end{lemma}

\begin{proof}
\omterm{def:b3:measure:outer}{बाह्य माप}: $\varnothing$ मनमाने छोटे अंतरालों से आच्छादित हो जाता है;
एकदिष्टता स्पष्ट है; उपयोज्यता: प्रत्येक $A_n$ के ऐसे आवरण दिए
हों जो निम्नतम से $\varepsilon 2^{-n}$ के भीतर हों, तो उनका सम्मिलन $\bigcup A_n$ को
कुल लंबाई $\leq \sum
\lambda^*(A_n) + \varepsilon$ के साथ आच्छादित कर देता है। स्थानांतरण
अपरिवर्तिता: आवरणों का स्थानांतरण कीजिए।

लंबाई: केवल $I = \intcc ab$ की स्थिति देखनी पर्याप्त है (अन्य प्रकार सिरों
से भिन्न होते हैं, जिनका \omterm{def:b3:measure:outer}{बाह्य माप} $0$ है: सूक्ष्म अंतरालों से
आच्छादित कीजिए; फिर $\intcc{a+\varepsilon}{b -
\varepsilon} \subseteq \intoo ab$ प्रकार की तुलनाओं से दबाइए)। $\lambda^*(\intcc ab) \leq b - a$:
$\intoo{a-\varepsilon}{b+\varepsilon}$ से आच्छादित कीजिए। विलोमतः मान लीजिए $\intcc
ab \subseteq \bigcup_n\intoo{a_n}{b_n}$: \emph{संहतता}
(बोरेल--लेबेग, \cref{thm:b3:topology:metriccompact}) से परिमित संख्या के अंतराल पर्याप्त हैं,
मान लीजिए $I_1, \dots, I_N$। हम $N$ पर आगमन से $\sum_{k\leq N}(b_k -
a_k) \geq b - a$ दिखाते हैं: $I_{k_1} \ni a$
चुनिए; यदि $b_{k_1} > b$ हो तो काम पूरा ($b_{k_1} - a_{k_1} > b - a$); अन्यथा रेखाखंड $\intcc{b_{k_1}}b$ शेष
$N - 1$ अंतरालों से आच्छादित है, और आगमन $\sum_{k \neq k_1}(b_k - a_k)
\geq b - b_{k_1}$ दे देता है, जबकि
$b_{k_1} - a_{k_1} > b_{k_1} - a$: अब जोड़ लीजिए।
\end{proof}

\begin{theorem}[लेबेग माप]\label{thm:b3:measure:lebesgue}
$\R$ का प्रत्येक बोरेल समुच्चय $\lambda^*$-मापनीय है। $\lambda^*$ का
$\sigma$-बीजगणित $\mathcal L = \mathcal M_{\lambda^*} \supseteq \mathcal B(\R)$ (\emph{लेबेग $\sigma$-बीजगणित}) पर प्रतिबंध
$\lambda$ $\mathcal B(\R)$ पर वह अद्वितीय \omterm{def:b3:measure:measure}{माप} है जो प्रत्येक अंतराल को उसकी
लंबाई सौंपता है; वह स्थानांतरण-अपरिवर्ती और $\sigma$-परिमित है।
\end{theorem}

\begin{proof}
\cref{thm:b3:measure:caratheodory} से केवल यह दिखाना पर्याप्त है कि प्रत्येक किरण $A = \intoo{-\infty}c$
$\lambda^*$-मापनीय है (किरणें $\mathcal B$ को जनित करती हैं, \cref{def:b3:measure:borel})। मान
लीजिए $\lambda^*(E) < \infty$ वाला $E \subseteq \R$ है और $\sum\ell(I_n) \leq \lambda^*(E) + \varepsilon$ वाला कोई आवरण $\bigcup I_n \supseteq E$। प्रत्येक
$I_n$ $\ell(I_n') + \ell(I_n'') = \ell(I_n)$ के साथ दो अंतरालों $I_n' = I_n \cap A$ और $I_n'' = I_n\setminus A$ में बँट जाता है
(किसी अंतराल में से किरण घटाने पर अंतराल ही बचता है); और $I_n'$
$E \cap A$ को तथा $I_n''$ $E \setminus A$ को आच्छादित करते हैं (परिभाषा के भीतर
रहने के लिए प्रत्येक को लंबाई $\ell +
\varepsilon2^{-n}$ के विवृत अंतराल तक बढ़ा
दीजिए), अतः
\[
\lambda^*(E\cap A) + \lambda^*(E\setminus A)
\leq \sum_n\bigl(\ell(I_n') + \ell(I_n'')\bigr) + 2\varepsilon
\leq \lambda^*(E) + 3\varepsilon .
\]
अद्वितीयता: अंतरालों की $\pi$-प्रणाली $\intoc ab$ पर लंबाई से सहमत दो
\omterm{def:b3:measure:measure}{माप} (जो उन पर परिमित हैं) $P_k = \intoc{-k}k$ के साथ \cref{thm:b3:measure:uniqueness} से $\sigma(\text{अंतराल}) = \mathcal B$ पर सहमत हो
जाते हैं। $\sigma$-परिमितता: $\R = \bigcup(-k, k]$।
\end{proof}

\begin{theorem}[नियमितता]\label{thm:b3:measure:regularity}
प्रत्येक $A \in \mathcal L$ के लिए:\index{लेबेग माप की नियमितता}
\[
\lambda(A) = \inf\{\lambda(U) : U \supseteq A \text{ विवृत}\}
= \sup\{\lambda(K) : K \subseteq A \text{ संहत}\}.
\]
\end{theorem}

\begin{proof}
\emph{बाह्य}: $\sum\ell(I_n) \leq
\lambda(A) + \varepsilon$ वाला कोई आवरण $\bigcup I_n$ $\lambda(U) \leq \lambda(A) + \varepsilon$ वाला \omterm{def:b3:topology:topology}{विवृत समुच्चय}
$U \supseteq A$ है (उपयोज्यता); और यदि $\lambda(A) = \infty$ हो तो कथन तुच्छ है।
\emph{आंतरिक}: पहले मान लीजिए $A$ परिबद्ध है, $A \subseteq [-M, M]$। $\lambda(U) \leq
\lambda([-M,M]\setminus A) + \varepsilon$ वाला
कोई विवृत $U \supseteq ([-M,M]\setminus A)$ चुनिए; तब $K = [-M,
M]\setminus U$ \omterm{def:b3:topology:compact}{संहत} है, $K \subseteq A$, और
\[
\lambda(K) = \lambda([-M,M]) - \lambda([-M,M]\cap U)
\geq \lambda([-M,M]) - \bigl(\lambda([-M,M]) - \lambda(A) +
\varepsilon\bigr) = \lambda(A) - \varepsilon .
\]
व्यापक $A$ के लिए: $\lambda(A) = \lim_M\lambda(A\cap[-M,M])$ (नीचे से \omterm{def:b3:topology:continuity}{सांतत्य}) और भीतर परिबद्ध
स्थिति लगाइए।
\end{proof}

\begin{example}\label{ex:b3:measure:cantor}
कैंटर समुच्चय (\cref{exo:b3:topology:10}) के लिए $\lambda(C) = 0$: $C \subseteq C_n$ लंबाई $3^{-n}$ के $2^n$
अंतरालों का सम्मिलन है, अतः $\lambda(C) \leq (2/3)^n \to 0$। एक अगणनीय शून्य समुच्चय ---
गणनीयता \omterm{def:b3:measure:measure}{माप} को नहीं देखती। विलोमतः, \emph{मोटे कैंटर समुच्चय}
(\cref{exo:b3:measure:5}) कहीं सघन नहीं हैं पर उनका \omterm{def:b3:measure:measure}{माप} धनात्मक है: अर्थात्
\omterm{def:b3:topology:topology}{संस्थिति} भी \omterm{def:b3:measure:measure}{माप} को नहीं देखती। सप्ताहांत समस्या इस पारस्परिक क्रिया
को उसके चौंकाने वाले निष्कर्ष तक ले जाती है: ऐसे लेबेग-मापनीय
समुच्चय हैं जो बोरेल नहीं हैं।
\end{example}

\begin{theorem}[विटाली]\label{thm:b3:measure:vitali}
$\R$ के \emph{समस्त} उपसमुच्चयों पर ऐसा कोई \omterm{def:b3:measure:measure}{माप} नहीं है जो
स्थानांतरण-अपरिवर्ती हो और प्रत्येक अंतराल को उसकी लंबाई सौंपता
हो। विशेष रूप से $\mathcal L \neq \mathcal P(\R)$: अमापनीय समुच्चय विद्यमान
हैं।\index{विटाली का अमापनीय समुच्चय}

\end{theorem}

\begin{proof}
मान लीजिए $\mu$ ऐसा कोई \omterm{def:b3:measure:measure}{माप} है। $\intcc01$ पर तुल्यता $x \sim y \iff x - y \in \Q$ पर विचार
कीजिए; \emph{चयन स्वयंसिद्ध} से प्रत्येक वर्ग में से $\intcc01$ का एक
प्रतिनिधि चुनिए: इससे एक समुच्चय $V$ बनता है। $q \in \Q\cap\intcc{-1}1$ के लिए
स्थानांतरण $V + q$ जोड़े-जोड़े में असंयुक्त हैं ($V$ के दो बिंदु
जिनका अंतर परिमेय हो, तुल्य होते हुए भी भिन्न प्रतिनिधि होते) और
\[
\intcc01 \subseteq \bigsqcup_{q \in \Q\cap\intcc{-1}1}(V + q)
\subseteq \intcc{-1}2 :
\]
जिसमें पहला अंतर्वेशन इसलिए कि प्रत्येक $x \in \intcc01$ अपने प्रतिनिधि से
किसी परिमेय $v$ जितना ही भिन्न होता है, $q = x - v \in
\intcc{-1}1$। एकदिष्टता और
$\sigma$-योज्यता से
\[
1 \leq \sum_{q}\mu(V + q) \leq 3,
\qquad\text{जहाँ } \mu(V + q) = \mu(V) \text{ प्रत्येक } q \text{ के लिए} .
\]
अचर $\mu(V)$ का अनंत योग $0$ या $\infty$ होता है: दोनों परिबंध एक
साथ लागू नहीं हो सकते। अतः ऐसा कोई $\mu$ विद्यमान नहीं है ---
और $V
\notin \mathcal L$, क्योंकि $\mathcal L$ पर $\lambda$ में प्रयुक्त सभी गुण मौजूद हैं।
\end{proof}

\begin{remark}\label{rem:b3:measure:banachtarski}
$\R^3$ में यह विफलता और भी नाटकीय है: बानाख--टार्स्की विरोधाभास
किसी गोले को पाँच टुकड़ों में विभाजित कर देता है जो घूर्णनों और
स्थानांतरणों से जुड़कर उसी त्रिज्या के \emph{दो} गोले बना देते हैं
--- अतः $\R^3$ के समस्त उपसमुच्चयों पर घूर्णन-अपरिवर्ती परिमित-योज्य
आयतन तक विद्यमान नहीं है। टुकड़े, निस्संदेह, अमापनीय हैं। मापनीयता
कोई नौकरशाही सावधानी नहीं है; वह संगति की सीमा है।
\end{remark}

\begin{method}\label{met:b3:measure:goodsets}
\emph{सुसमुच्चय सिद्धांत}: यह सिद्ध करने के लिए कि $\sigma(\mathcal E)$ के समस्त
समुच्चयों में कोई गुण है, दिखाइए कि अच्छे समुच्चय $\mathcal E$ को
अंतर्विष्ट करने वाला $\sigma$-बीजगणित (अथवा $\lambda$-प्रणाली बनाते हैं,
यदि गुण माप-सैद्धांतिक हो और $\mathcal E$ एक $\pi$-प्रणाली हो --- तब
डिन्किन)। इस अध्याय की और अगले अध्याय की लगभग हर उपपत्ति इसी का
दृष्टांत है। दो मापों को बराबर सिद्ध करने के लिए: उन्हें किसी जनक
$\pi$-प्रणाली पर जाँचिए और $\sigma$-परिमितता भी (\cref{thm:b3:measure:uniqueness})। कोई \omterm{def:b3:measure:measure}{माप}
बनाने के लिए: आवरणों से एक \omterm{def:b3:measure:outer}{बाह्य माप} बनाइए और कैराथियोडोरी उद्धृत
कीजिए।
\end{method}

\section{अभ्यास}

\begin{exercise}[$\star$]\label{exo:b3:measure:1}
(क) दर्शाइए कि $\{A \subseteq X : A$ अथवा $X\setminus A$ गणनीय है$\}$ एक $\sigma$-बीजगणित
है: वही जो एकल-अवयवी समुच्चयों से जनित होता है।
(ख) दर्शाइए कि $\mathcal B(\R)$ इनमें से प्रत्येक से जनित होता है: विवृत
अंतराल; संवृत अंतराल; किरणें $\intoc{-\infty}a$; और $a \in \Q$ वाली किरणें।
(ग) क्या \emph{अंतरालों के परिमित असंयुक्त सम्मिलनों} $\intoc ab$ का
कुल एक $\sigma$-बीजगणित है? क्या वह एक बीजगणित है (पूरक और परिमित
सम्मिलन के अंतर्गत स्थिर)?
\end{exercise}

\begin{exercise}[$\star$]\label{exo:b3:measure:2}
(क) किसी परिमित \omterm{def:b3:measure:measure}{माप} के लिए समावेशन--अपवर्जन सिद्ध कीजिए: $\mu(A\cup
B) = \mu(A) + \mu(B) - \mu(A\cap B)$,
और तीन समुच्चयों वाला रूप भी।
(ख) ऐसा उदाहरण दीजिए जो दर्शाए कि परिमितता परिकल्पना के बिना ऊपर
से \omterm{def:b3:topology:continuity}{सांतत्य} (\cref{prop:b3:measure:basics}(घ)) विफल हो जाती है।
(ग) दर्शाइए कि किसी गणनीय समुच्चय का \omterm{def:b3:measure:lebesgueouter}{लेबेग माप} शून्य होता है।
इससे $\lambda(\Q) = 0$ और $\lambda(\intcc01\setminus\Q) =
1$ निष्कर्ष निकालिए।
\end{exercise}

\begin{exercise}[$\star\star$]\label{exo:b3:measure:3}
मान लीजिए $\mu, \nu$ $\mathcal B(\R)$ पर प्रायिकता \omterm{def:b3:measure:measure}{माप} हैं जिनके लिए सभी $t \in \R$ पर
$\mu(\intoc{-\infty}t) = \nu(\intoc{-\infty}t)$। दर्शाइए कि $\mu = \nu$। (इससे \emph{बंटन फलन} $F(t) = \mu(\intoc{-\infty}t)$ एक \omterm{def:b3:complete:complete}{पूर्ण}
अपरिवर्ती बन जाता है --- यही \cref{ch:b3:probability} की नींव है।)
\end{exercise}

\begin{exercise}[$\star\star$]\label{exo:b3:measure:4}
(बोरेल--कैंटेली, \omterm{def:b3:measure:measure}{माप} रूप) मान लीजिए $(A_n)$ मापनीय हैं जिनके लिए
$\sum_n\mu(A_n) < \infty$, और $\limsup A_n =
\bigcap_N\bigcup_{n\geq N}A_n$ (वे बिंदु जो अनंत संख्या के $A_n$ में हैं)।
दर्शाइए कि $\mu(\limsup A_n) = 0$। अनुप्रयोग: लगभग प्रत्येक $x \in \intcc01$ के लिए, $\Q\cap\intcc01$
की किसी गणना के $n$-वें परिमेय $p/q_n$ हेतु केवल परिमित संख्या
के $n$ $\abs{x - p/q_n} \leq 4^{-n}$ को संतुष्ट करते हैं।
\end{exercise}

\begin{exercise}[$\star\star$]\label{exo:b3:measure:5}
(मोटा कैंटर समुच्चय) $\intcc01$ पर कैंटर रचना दोहराइए, पर चरण $n$ पर
$2^{n-1}$ अंतरालों में से प्रत्येक से केवल लंबाई $4^{-n}$ का
\emph{केंद्रित} विवृत अंतराल हटाइए। दर्शाइए कि प्राप्त $K = \bigcap K_n$ \omterm{def:b3:topology:compact}{संहत}
है, उसका \omterm{def:b3:topology:interior}{अंतःभाग} रिक्त है (कोई अंतराल नहीं बचता), और
\[
\lambda(K) = 1 - \sum_{n\geq1}2^{n-1}4^{-n} = \tfrac12 :
\]
अर्थात् \omterm{def:b3:measure:measure}{माप} $\frac12$ वाला एक कहीं सघन नहीं समुच्चय। इससे $\intcc01$ का
पूरा \omterm{def:b3:measure:measure}{माप} $1$ वाला एक \emph{\omterm{rem:b3:complete:meagre}{क्षीण}} उपसमुच्चय निकालिए, और \omterm{def:b3:measure:measure}{माप}
$< \varepsilon$ वाला एक विवृत सघन उपसमुच्चय भी।
\end{exercise}

\begin{exercise}[$\star\star$]\label{exo:b3:measure:6}
मान लीजिए $\mu$ $\mathcal B(\R)$ पर स्थानांतरणों के अंतर्गत अपरिवर्ती ऐसा
\omterm{def:b3:measure:measure}{माप} है जिसके लिए $c = \mu(\intoc01) < \infty$। दर्शाइए कि $\mathcal B(\R)$ पर $\mu =
c\,\lambda$। \emph{($\intoc01$
को $2^n$ स्थानांतरणों में बाँटकर द्विआधारी अंतरालों पर $\mu$
परिकलित कीजिए, फिर \cref{thm:b3:measure:uniqueness} का आह्वान कीजिए।)}
\end{exercise}

\begin{exercise}[$\star\star$]\label{exo:b3:measure:7}
(आसन्नन) मान लीजिए $A \in \mathcal L$ है, $\lambda(A) <
\infty$, और $\varepsilon > 0$। दर्शाइए कि
अंतरालों का ऐसा \emph{परिमित} सम्मिलन $B$ है जिसके लिए $\lambda(A\,\triangle\,B) <
\varepsilon$
($\triangle$ = सममित अंतर)। \emph{(नियमितता: $K \subseteq A \subseteq U$ को दबाइए और विवृत
$U$ की अंतरालों के गणनीय सम्मिलन वाली संरचना तथा $K$ की
संहतता का उपयोग कीजिए।)}
\end{exercise}

\begin{exercise}[$\star\star\star$]\label{exo:b3:measure:8}
(श्टाइनहाउस) मान लीजिए $A \in \mathcal L$ है और $\lambda(A) > 0$। दर्शाइए कि $A - A = \{x - y : x, y \in A\}$ $0$
के चारों ओर एक अंतराल अंतर्विष्ट करता है। \emph{($\lambda(A) < \infty$ तक सीमित
कीजिए; \cref{exo:b3:measure:7} जैसी नियमितता से $\lambda(A \cap I) > \frac34\ell(I)$ वाला कोई अंतराल $I$
ढूँढ़िए; फिर $\abs t <
\frac12\ell(I)$ के लिए समुच्चय $A\cap I$ और $(A\cap I) + t$ दोनों लंबाई
$\frac32\ell(I)$ के किसी अंतराल में बैठते हैं और उनका कुल \omterm{def:b3:measure:measure}{माप} $> \frac32\ell(I)$ है: अतः
उन्हें प्रतिच्छेद करना ही होगा।)}
\end{exercise}

\begin{exercise}[$\star\star\star$]\label{exo:b3:measure:9}
दर्शाइए कि $\lambda(A) > 0$ वाला प्रत्येक $A \in \mathcal L$ किसी अमापनीय उपसमुच्चय को
अंतर्विष्ट करता है। \emph{($A$ को विटाली के समुच्चय के
स्थानांतरणों $V + q$ से प्रतिच्छेद कीजिए: यदि सभी $A \cap (V+q)$ मापनीय
होते, तो \cref{thm:b3:measure:vitali} के तर्क से प्रत्येक शून्य होता --- श्टाइनहाउस
(\cref{exo:b3:measure:8}) सहायक है: $V + q$ के भीतर धनात्मक \omterm{def:b3:measure:measure}{माप} वाला कोई मापनीय
समुच्चय $(V+q) - (V+q) \supseteq$ को एक अंतराल बना देता, जो इस बात का खंडन है कि यह
अंतर समुच्चय $\Q$ से केवल $0$ पर मिलता है; उपयोज्यता से
निष्कर्ष निकालिए।)}
\end{exercise}

\begin{exercise}[$\star\star$]\label{exo:b3:measure:10}
दर्शाइए कि $\lambda^*(A) < \infty$ वाला $A \subseteq \R$ लेबेग-मापनीय है तभी और केवल तभी जब
प्रत्येक $\varepsilon > 0$ के लिए $\lambda^*(U \setminus A) <
\varepsilon$ वाला कोई विवृत $U \supseteq A$ हो, और तभी और
केवल तभी जब $\lambda^*(G\setminus A) = 0$ वाला कोई $G_\delta$ समुच्चय $G \supseteq A$ हो। (अर्थात् लेबेग
समुच्चय शून्य समुच्चयों तक बोरेल समुच्चय ही हैं।)
\end{exercise}

\begin{exercise}[$\star\star$]\label{exo:b3:measure:11}
(एकदिष्ट सीमाओं के अनुदिश \omterm{def:b3:topology:continuity}{सांतत्य}, और उसकी तीक्ष्णता)
(क) दर्शाइए कि मापनीय समुच्चयों के लिए $\mu(\liminf A_n) \leq
\liminf\mu(A_n)$ (समुच्चयों के लिए
फतू), और यह कि यदि $\mu\bigl(\bigcup A_n\bigr) < \infty$ हो तो $\limsup\mu(A_n) \leq \mu(\limsup A_n)$ भी।
(ख) $\R$ पर \omterm{def:b3:measure:lebesgueouter}{लेबेग माप} के लिए ऐसा अनुक्रम प्रस्तुत कीजिए जिसके
लिए सभी $n$ पर $\mu(A_n) = 1$ हो फिर भी $\mu(\limsup A_n) = 0$ हो: दूसरी असमिका की
परिमितता परिकल्पना केवल सजावट नहीं है।
(ग) निष्कर्ष निकालिए: यदि $\sum\mu(A_n) < \infty$ हो तो $\mu(\limsup A_n) = 0$ (फिर बोरेल--कैंटेली),
और यदि $A_n$ बढ़ते या घटते हों (घटती स्थिति में $\mu(A_1) < \infty$ के साथ) तो
$\mu(\lim A_n) = \lim\mu(A_n)$।
\end{exercise}

\begin{exercise}[$\star\star\star$]\label{exo:b3:measure:12}
(एगोरोव की प्रमेय) मान लीजिए $\mu(X) < \infty$ और बिंदुशः $f_n \to f$, और सभी फलन
मापनीय (वास्तविक-मान वाले) हैं। $k, N \geq 1$ के लिए रखिए
\[
E_{k,N} = \bigcap_{n \geq N}\Bigl\{x : \abs{f_n(x) - f(x)}
\leq \tfrac1k\Bigr\} .
\]
(क) दर्शाइए कि स्थिर $k$ के लिए $N \to
\infty$ पर $E_{k,N} \nearrow X$, और $\mu(X \setminus E_{k,N_k})
\leq \varepsilon2^{-k}$ वाला
$N_k$ निष्कर्ष के रूप में निकालिए।
(ख) \emph{एगोरोव की प्रमेय} निष्कर्ष के रूप में निकालिए: प्रत्येक
$\varepsilon
> 0$ के लिए $\mu(X\setminus A) \leq
\varepsilon$ वाला ऐसा मापनीय $A$ है कि \emph{$A$ पर
एकसमान रूप से} $f_n \to f$ --- अर्थात् बिंदुशः अभिसरण मनमाने छोटे
समुच्चय के बाहर एकसमान अभिसरण होता है।
(ग) दर्शाइए कि $(\R, \lambda)$ पर यह प्रमेय विफल हो जाती है: चलती हुई
उभारें $f_n = \mathbf 1_{\intcc n{n+1}}$ बिंदुशः $0$ पर अभिसरित होती हैं पर परिमित \omterm{def:b3:measure:measure}{माप} वाले
किसी भी समुच्चय के पूरक पर एकसमान रूप से नहीं। (क) में $\mu(X) < \infty$ का
उपयोग कहाँ हुआ?
\end{exercise}

\section{समस्या: कैंटर--विटाली सीढ़ी और एक ऐसा मापनीय समुच्चय जो
बोरेल नहीं है}

\begin{omfigure}
\begin{tikzpicture}[scale=5.2]
  \draw[->] (0,0) -- (1.06,0) node[right, font=\small] {$x$};
  \draw[->] (0,0) -- (0,1.06);
  \draw[black!30] (1,0) -- (1,1) -- (0,1);
  \draw[omThm, thick]
    (0,0)
    -- (0.037,0.125) -- (0.074,0.125)
    -- (0.111,0.25)  -- (0.222,0.25)
    -- (0.259,0.375) -- (0.296,0.375)
    -- (0.333,0.5)   -- (0.667,0.5)
    -- (0.704,0.625) -- (0.741,0.625)
    -- (0.778,0.75)  -- (0.889,0.75)
    -- (0.926,0.875) -- (0.963,0.875)
    -- (1,1);
  \node[font=\small, omThm] at (0.32,0.86)
    {$c$ (गहराई-3 आसन्नन)};
\end{tikzpicture}

{\small कैंटर--विटाली सीढ़ी: कैंटर समुच्चय की प्रत्येक रिक्ति पर
अचर, फिर भी $0$ से $1$ तक संततता से चढ़ती हुई। उसका अवकलज
लगभग सर्वत्र लुप्त हो जाता है --- सारी चढ़ाई एक शून्य समुच्चय पर
होती है।}
\end{omfigure}

\begin{problem}[{सप्ताहांत समस्या --- शैतान की सीढ़ी, और $\mathcal B(\R) \subsetneq \mathcal L$}]
\label{pb:b3:measure:1}
हम \emph{कैंटर--विटाली फलन} (शैतान की सीढ़ी) की रचना करते हैं,
उससे \omterm{def:b3:measure:measure}{माप} का विकृत परिवहन करते हैं, और एक ऐसी प्रमेय पर पहुँचते
हैं जो किसी भी मुलायम तर्क से नहीं मिलती: ऐसे लेबेग-मापनीय
समुच्चय विद्यमान हैं जो बोरेल नहीं हैं। संकेतन: $C$ कैंटर
समुच्चय है, $C_n$ उसका $n$-वाँ चरण (लंबाई $3^{-n}$ के $2^n$
अंतराल), और प्रत्येक $x \in C$ के त्र्याधारी अंक $x = \sum 2b_n3^{-n}$, $b_n \in \{0,1\}$
(\cref{exo:b3:topology:10}) हैं।

\medskip
\noindent\textbf{भाग I --- सीढ़ी।} $c_0(x) = x$ परिभाषित कीजिए और $c_n$
से $c_{n+1}$ इस प्रकार:
\[
c_{n+1}(x) = \begin{cases}
\tfrac12\,c_n(3x) & 0 \leq x \leq \tfrac13,\\[2pt]
\tfrac12 & \tfrac13 \leq x \leq \tfrac23,\\[2pt]
\tfrac12 + \tfrac12\,c_n(3x - 2) & \tfrac23 \leq x \leq 1.
\end{cases}
\]
\begin{enumerate}
  \item दर्शाइए कि प्रत्येक $c_n$ \omterm{def:b3:topology:continuity}{संतत}, अनवरोही है, जिसमें
        $c_n(0) = 0$, $c_n(1) = 1$, और यह कि $\norm{c_{n+1} - c_n}_\infty \leq
        \tfrac12\norm{c_n - c_{n-1}}_\infty$।
  \item निष्कर्ष निकालिए कि $(c_n)$ एकसमान रूप से किसी \omterm{def:b3:topology:continuity}{संतत}
        अनवरोही $c$ पर अभिसरित होता है जिसके लिए $c(0) = 0$, $c(1) = 1$
        (\emph{वही} कैंटर--विटाली फलन), और जो ऊपर के $c_{n+1}$ जैसे
        ही स्वसमरूप संबंध संतुष्ट करता है।
  \item दर्शाइए कि $c$ $\intcc01\setminus C$ के प्रत्येक \omterm{def:b3:topology:components}{संबद्ध घटक} पर अचर है,
        और यह कि $x = \sum_n
        2b_n3^{-n} \in C$ के लिए $c(x) = \sum_n b_n2^{-n}$: अर्थात् सीढ़ी कैंटर अंकों को
        द्व्याधारी रूप में पढ़ती है (\cref{pb:b3:topology:1} का फलन $g$, जिसे
        एकदिष्ट और वैश्विक बना दिया गया है)।
  \item निष्कर्ष निकालिए कि $c$ $\intcc01\setminus C$ के प्रत्येक बिंदु पर
        $c' = 0$ के साथ अवकलनीय है: अर्थात् \emph{$\lambda$-लगभग सर्वत्र}
        $c' = 0$ (\cref{ex:b3:measure:cantor})। निष्कर्ष निकालिए कि $c$ के लिए कलन की
        मूल प्रमेय विफल हो जाती है:
        \[
        c(1) - c(0) = 1 \neq 0 = \int_0^1 c'(t)\,\dd t
        \]
        (समाकल उस पूरे-माप वाले समुच्चय पर है जहाँ $c' = 0$; \cref{ch:b3:lebesgue}
        की प्रत्याशा में, शून्य समुच्चय समाकलों को प्रभावित नहीं
        करते)। $\mathcal C^1$ की मूल प्रमेय की कौन-सी परिकल्पना का उल्लंघन
        होता है?
  \item दर्शाइए कि $c(C) = \intcc01$: शून्य समुच्चय $C$ पूरे \omterm{def:b3:measure:measure}{माप} वाले एक
        समुच्चय \emph{पर} प्रतिचित्रित हो जाता है।
\end{enumerate}

\medskip
\noindent\textbf{भाग II --- टेढ़ी \omterm{def:b3:topology:continuity}{समस्थितिकता}।} मान लीजिए $h(x) = \frac{c(x) + x}{2}$।
\begin{enumerate}[resume]
  \item दर्शाइए कि $h \colon \intcc01 \to \intcc01$ एक \omterm{def:b3:topology:continuity}{समस्थितिकता} है (कड़ाई से बढ़ता हुआ,
        \omterm{def:b3:topology:continuity}{संतत}, आच्छादक)।
  \item दर्शाइए कि $\lambda\bigl(h(\intcc01\setminus C)\bigr) =
        \tfrac12$: लंबाई $\ell$ की प्रत्येक रिक्ति पर $h$
        ढाल $\tfrac12$ के आफ़ीन प्रतिचित्रण की तरह क्रिया करता है, और
        रिक्तियों की कुल लंबाई $1$ है।
  \item $\lambda\bigl(h(C)\bigr) = \tfrac12$ निष्कर्ष के रूप में निकालिए: किसी शून्य समुच्चय के
        समस्थितिक प्रतिबिंब का \omterm{def:b3:measure:measure}{माप} धनात्मक हो सकता है। (यह
        ``आकार'' के बारे में भोली अंतर्दृष्टि का खंडन कहाँ करता
        है?)
\end{enumerate}

\medskip
\noindent\textbf{भाग III --- एक ऐसा मापनीय समुच्चय जो बोरेल नहीं
है।}
\begin{enumerate}[resume]
  \item \cref{exo:b3:measure:9} से कोई अमापनीय $W
        \subseteq h(C)$ चुनिए। दर्शाइए कि $Z = h^{-1}(W) \subseteq C$
        लेबेग-मापनीय है। \emph{(वह किसी शून्य समुच्चय का
        उपसमुच्चय है; \omterm{def:b3:complete:complete}{पूर्णता}, \cref{thm:b3:measure:caratheodory}।)}
  \item दर्शाइए कि किसी \omterm{def:b3:topology:continuity}{संतत प्रतिचित्रण} के अंतर्गत किसी बोरेल
        समुच्चय का पूर्वप्रतिबिंब बोरेल होता है। \emph{(सुसमुच्चय
        सिद्धांत: $\{B : h^{-1}(B) \in \mathcal B\}$ विवृत समुच्चयों को अंतर्विष्ट करने वाला
        $\sigma$-बीजगणित है --- प्रतिचित्रण की दिशा का ध्यान
        रखिए।)}
  \item निष्कर्ष निकालिए कि $Z$ बोरेल \emph{नहीं} है: यदि वह
        होता, तो $W = (h^{-1})^{-1}(Z)$ बोरेल होता (\omterm{def:b3:topology:continuity}{संतत} $h^{-1}$ पर प्रश्न 10
        लगाइए), और इसलिए मापनीय --- विरोधाभास। अतः
        \[
        \boxed{\ \mathcal B(\R) \subsetneq \mathcal L\ }
        \]
        और \omterm{def:b3:complete:complete}{पूर्णता} वस्तुतः बोरेल संसार को बड़ा कर देती है।
  \item ऐसा लेबेग-मापनीय फलन $g$ और ऐसा \omterm{def:b3:topology:continuity}{संतत} फलन $\varphi$
        प्रस्तुत कीजिए कि $g \circ
        \varphi$ लेबेग-मापनीय न हो: मापनीयता,
        \omterm{def:b3:topology:continuity}{सांतत्य} के विपरीत, संयोजित नहीं होती। \emph{(\cref{ch:b3:lebesgue} से
        मापनीय फलनों की परिभाषा की प्रत्याशा करते हुए $g =
        \mathbf 1_Z$ और
        $\varphi = h^{-1}$ लीजिए: बोरेल समुच्चयों के पूर्वप्रतिबिंब लेबेग
        समुच्चय होते हैं। लक्ष्य पर कौन-सा $\sigma$-बीजगणित प्रयुक्त
        है, इस विषय में कहाँ सावधानी बरतनी चाहिए?)}
\end{enumerate}

\medskip
\noindent\textbf{भाग IV --- उपसंहार।}
\begin{enumerate}[resume]
  \item निम्नलिखित वर्गों को कड़े अंतर्वेशन के अनुसार क्रमबद्ध
        कीजिए और प्रत्येक कड़ेपन को इस अध्याय और उसकी समस्या के
        किसी उदाहरण से उचित ठहराइए: गणनीय समुच्चय; बोरेल शून्य
        समुच्चय; लेबेग शून्य समुच्चय; बोरेल समुच्चय; लेबेग
        समुच्चय; मनमाने समुच्चय।
\end{enumerate}

\medskip
\noindent\textbf{भाग V --- कैंटर \omterm{def:b3:measure:measure}{माप}: शून्य समुच्चय पर
द्रव्यमान।} सीढ़ी एक उल्लेखनीय \omterm{def:b3:measure:measure}{माप} का बंटन फलन है, जिसे अब हम इसी
अध्याय के औजारों से बनाते हैं।
\begin{enumerate}[resume]
  \item (लेबेग--स्टीलत्येस, अस्तित्व) मान लीजिए $F\colon\R\to\R$ अनवरोही,
        \omterm{def:b3:topology:continuity}{संतत} और परिबद्ध है। अर्ध-विवृत अंतरालों पर $\rho\bigl(\intoc ab\bigr) = F(b) -
        F(a)$
        परिभाषित कीजिए और $A \subseteq \R$ के लिए
        \[
        \mu_F^*(A) = \inf\Bigl\{\sum_k\bigl(F(b_k) -
        F(a_k)\bigr) : A \subseteq
        \bigcup_k\intoc{a_k}{b_k}\Bigr\} .
        \]
        दर्शाइए कि $\mu_F^*$ एक \omterm{def:b3:measure:outer}{बाह्य माप} है और
        $\mu_F^*\bigl(\intoc ab\bigr) = F(b) - F(a)$
        \emph{(\cref{thm:b3:measure:lebesgue} के संहतता तर्क का अनुकरण कीजिए, प्रत्येक
        $\intoc{a_k}{b_k}$ को $\leq \varepsilon2^{-k}$ $F$-लागत पर किसी विवृत अंतराल तक बढ़ाते
        हुए --- $F$ का \omterm{def:b3:topology:continuity}{सांतत्य} कहाँ प्रयुक्त होता है?)}।
  \item दर्शाइए कि प्रत्येक बोरेल समुच्चय कैराथियोडोरी अर्थ में
        $\mu_F^*$-मापनीय है \emph{(लेबेग स्थिति की तरह, केवल
        अर्ध-रेखाओं की परीक्षा पर्याप्त है; \cref{thm:b3:measure:caratheodory} के अनुप्रयोग की
        उपपत्ति का अनुसरण कीजिए)}, अतः $\mathcal B(\R)$ पर प्रतिबंधित $\mu_F = \mu_F^*$
        $\mu_F(\intoc ab) = F(b) - F(a)$ वाला एक \omterm{def:b3:measure:measure}{माप} है: यही $F$ का \emph{लेबेग--स्टीलत्येस
        \omterm{def:b3:measure:measure}{माप}} है।
  \item इसे सीढ़ी पर लगाइए ($F = c$ को $\R_-$ पर $0$ से और $\intco1\infty$
        पर $1$ से विस्तारित कीजिए): इस प्रकार \emph{कैंटर \omterm{def:b3:measure:measure}{माप}}
        $\mu$ मिलता है। दर्शाइए कि $\mu(\R) = 1$, कि कैंटर समुच्चय की
        प्रत्येक रिक्ति $\mu$-शून्य है (वहाँ $c$ अचर है), और
        निष्कर्ष निकालिए
        \[
        \mu(C) = 1, \qquad \lambda(C) = 0 :
        \]
        अर्थात् $\mu$ और $\lambda$ असंयुक्त वाहकों ($C$ और उसके
        पूरक) पर रहते हैं। इस स्थिति वाले दो \omterm{def:b3:measure:measure}{माप} \emph{परस्पर
        एकल} कहलाते हैं, और $\mu \perp
        \lambda$ लिखा जाता है।
  \item दर्शाइए कि $\mu$ में कोई परमाणु नहीं है: प्रत्येक $x$
        के लिए $\mu(\{x\}) = 0$ \emph{($c$ का \omterm{def:b3:topology:continuity}{सांतत्य})}। किसी
        लेबेग-शून्य \omterm{def:b3:topology:compact}{संहत} समुच्चय पर धारण किया गया एक अपरमाणुक
        प्रायिकता \omterm{def:b3:measure:measure}{माप}: अब तक देखे गए मापों से इसकी तुलना कीजिए।
  \item (छद्म वेश में सिक्का उछाल) किसी शब्द $(\varepsilon_1,
        \dots, \varepsilon_m) \in \{0,1\}^m$ के लिए मान
        लीजिए $C_{\varepsilon}$ उन $x \in C$ का समुच्चय है जिनके त्र्याधारी अंक
        $i \leq m$ के लिए $b_i(x) = \varepsilon_i$ को संतुष्ट करते हैं (अर्थात् गहराई
        $m$ के $2^m$ कैंटर टुकड़ों में से एक)। दर्शाइए कि
        $\mu(C_\varepsilon) = 2^{-m}$ \emph{(उस टुकड़े के पार सीढ़ी $2^{-m}$ चढ़ती है:
        भाग I, प्रश्न 3 का उपयोग कीजिए)}। कैंटर \omterm{def:b3:measure:measure}{माप} त्र्याधारी
        रूप में पढ़े गए न्यायसंगत सिक्कों की अनंत उछालों का नियम
        है --- \cref{ch:b3:probability} इसे यथार्थ बना देगा।
  \item स्वसमरूपता सिद्ध कीजिए: प्रत्येक बोरेल $A$ के लिए
        \[
        \mu(A) = \tfrac12\,\mu(3A) +
        \tfrac12\,\mu(3A - 2) ,
        \]
        जहाँ $3A - 2 = \{3x - 2 : x \in A\}$ \emph{($c$ के स्वसमरूप संबंधों के माध्यम से
        जनक अंतरालों $\intoc ab$ पर इसे जाँचिए, फिर अद्वितीयता \cref{thm:b3:measure:uniqueness}
        का आह्वान कीजिए)}।
  \item दर्शाइए कि परावर्तन $s(x) = 1 - x$ $\mu$ को सुरक्षित रखता है:
        $\mu(s(A)) = \mu(A)$ \emph{($c(1 - x) = 1 -
        c(x)$ के माध्यम से, जो रचना की सममिति से
        निकलता है --- उसे सिद्ध कीजिए)}।
  \item $\mu$ के पहले दो आघूर्ण परिकलित कीजिए, अर्थात् नियम
        $\mu$ वाले किसी यादृच्छिक बिंदु $X$ के (समाकलों को
        \cref{ch:b3:lebesgue} की प्रत्याशा में गहराई-$m$ टुकड़ों पर योगों की
        सीमा के रूप में लिया जा सकता है): सममिति से $\int x\,\dd\mu = \frac12$, और
        स्वसमरूपता से
        \[
        \int x^2\,\dd\mu = \frac38,
        \qquad\text{अतः}\qquad
        \operatorname{Var}(X) = \frac18 .
        \]
        $\intcc01$ पर एकसमान नियम से तुलना कीजिए (प्रसरण $\frac1{12}$): कैंटर
        द्रव्यमान, किनारों की ओर धकेला हुआ, \emph{अधिक} फैलता है।
  \item दर्शाइए कि $\mu$ का (संस्थितिक) \emph{आलंब} --- अर्थात्
        पूरे \omterm{def:b3:measure:measure}{माप} वाला सबसे छोटा संवृत समुच्चय --- ठीक $C$ है।
  \item (संश्लेषण) सीढ़ी $c$ \omterm{def:b3:topology:continuity}{संतत} और अनवरोही है, फिर भी कलन की
        मूल प्रमेय को संतुष्ट नहीं करती (भाग I); और \omterm{def:b3:measure:measure}{माप} $\mu_c$ एक
        प्रायिकता है, अपरमाणुक है, तथा $\lambda$ के सापेक्ष एकल है।
        एक छोटे अनुच्छेद में समझाइए कि ये एक ही परिघटना के दो
        चेहरे कैसे हैं, और व्यापक सार कहिए: अनवरोही फलन मापों के
        अनुरूप होते हैं ($F \leftrightarrow \mu_F$), लगभग सर्वत्र अवकलनीयता
        ``निरपेक्षतः \omterm{def:b3:topology:continuity}{संतत} भाग'' के अनुरूप होती है, और $c$ वह
        मानक साक्षी है कि कोई \omterm{def:b3:topology:continuity}{संतत} $F$ बिलकुल \emph{कोई}
        निरपेक्षतः \omterm{def:b3:topology:continuity}{संतत} भाग धारण न करे, यह भी संभव है।
  \item (\omterm{def:b3:topology:continuity}{सांतत्य} का तीक्ष्ण मापांक) मान लीजिए $s =
        \frac{\ln 2}{\ln 3}$। दर्शाइए कि
        $c$ घातांक $s$ का होल्डर \omterm{def:b3:topology:continuity}{संतत} है:
        \[
        \abs{c(x) - c(y)} \leq 4\,\abs{x - y}^{s}
        \qquad (x, y \in \intcc01),
        \]
        और यह कि $t > s$ वाला कोई घातांक स्थानीय रूप से भी काम नहीं
        कर सकता। माप-सैद्धांतिक रूप निकालिए: प्रत्येक $x$ और
        प्रत्येक $r \in \intoc01$ के लिए
        \[
        \mu\bigl(\intcc{x - r}{x + r}\bigr) \leq 8\,r^{s} .
        \]
        \emph{(गहराई-$m$ की त्र्याधारी जालिका की तुलना $\abs{x - y}$ के
        मापक्रम से कीजिए; प्रश्न 18 प्रत्येक टुकड़े के पार
        चढ़ाई दे देता है। घातांक $s$ $C$ की \omterm{def:b3:topology:hausdorff}{हाउसडॉर्फ} विमा है,
        जैसा आगे के पाठ्यक्रम कहेंगे।)}
  \item (स्वसमरूपता $\mu$ को अभिलक्षित कर देती है) प्रश्न 19 का
        विलोम सिद्ध कीजिए: यदि $\nu$ $\mathcal B(\R)$ पर $\intcc01$ द्वारा धारण
        किया गया ऐसा प्रायिकता \omterm{def:b3:measure:measure}{माप} हो जो
        \[
        \nu(A) = \tfrac12\,\nu(3A) + \tfrac12\,\nu(3A - 2)
        \qquad (A \in \mathcal B(\R)),
        \]
        को संतुष्ट करे, तो $\nu = \mu$। \emph{($\nu$ को गहराई-$m$ के
        $2^m$ कैंटर टुकड़ों पर फैलाने के लिए संबंध को $m$ बार
        दोहराइए, $\intoc ab$ के भीतर के टुकड़ों की गणना के सापेक्ष
        $\nu(\intoc ab)$ का आकलन कीजिए, और $m \to \infty$ जाने दीजिए; \cref{thm:b3:measure:uniqueness} से पूरा
        कीजिए।)}
\end{enumerate}
\end{problem}
